\documentclass[12pt]{article}

\usepackage{amssymb, amsmath, amsfonts, amsthm, graphicx, tcolorbox}
\usepackage{tikz-cd}

\parskip = .2in
\parindent  = 0in

\pagestyle{empty}


%==========
%==========

\begin{document}

\begin{titlepage}
Math 174, Fall 2026 \\ 
Homework 2 \\

\vfill

% If you are typesetting your work, just remove the entire \underline bit and put your name there instead. 
Name: Jayden Thadani


\end{titlepage}

%==========
\begin{tcolorbox}[title=Exercise 1]
Find at least three pairwise inequivalent degree 2 representations of $S_2$. 
\end{tcolorbox}

\bigskip

Write $S_{2} \cong \mathbb{Z}^{\times}$ multiplicatively for convenience. Consider the following three representations.

\begin{enumerate}
	\item The trivial representation $\rho_{1} : S_{2} \to \mathrm{GL}_{2}(\mathbb{C})$ by $g \mapsto I = \begin{pmatrix}
			1 & 0 \\
			0 & 1
	\end{pmatrix}$ for each $g \in S_{2}$.

	\item The permutation representation $\rho_{2}$ by $-1 \mapsto \begin{pmatrix}
		0 & 1 \\
		1 & 0
	\end{pmatrix}$. This completely describes the representation since $-1$ generates $\mathbb{Z}^{\times}$. 

	Note, in class we showed that $\rho_{2} \cong \rho_{2}'$ where $\rho_{2}' : -1 \mapsto \begin{pmatrix}
		1 & 0 \\
		0 & -1
	\end{pmatrix}$.

\item The representation $\rho_{3}$ by $-1 \mapsto \begin{pmatrix}
		-1 & 0 \\
		0 & -1
\end{pmatrix}$.
\end{enumerate}

$\rho_{1} \not \cong \rho_{2}$ because any change of basis must give $T \rho_{1}(-1) T^{-1} = T I T ^{-1} = I$ and $\rho_{2}(-1) \neq I$. $\rho_{1} \not \cong \rho_{3}$ for the same reason.

$\rho_{2} \neq \rho_{3}$ because they have different decompositions into irreducible pieces. In particular, there are two (irreducible) $1$-dimensional representations $\pi_{1}, \pi_{2} : \mathbb{Z}^{\times} \to \mathbb{C}^{\times}$. $\pi_{1}$ is trivial and $\pi_{2}$ is given by $-1 \mapsto -1$. As previously, $\pi_{1} \not \cong \pi_{2}$. While $\rho_{2} \cong \rho_{2}' \cong \pi_{1} \oplus \pi_{2}$, we have $\rho_{3} \cong \pi_{2} \oplus \pi_{2}$.

\newpage

%==========
\begin{tcolorbox}[title=Exercise 2]
Let $G$ be a cyclic group of order $n$. Suppose $g \in G$ is a generator for $G$ so that $G = \{ 1, g, g^2, \dots, g^{n-1} \}$. How many pairwise inequivalent degree 1 representations does $G$ have? 

\bigskip

[Hint:  Use the primitive $n$-th root of unity $\omega = e^{2 \pi i / n}$ in your solution.] 
\end{tcolorbox}

\bigskip

Suppose $\rho : G \to \mathrm{GL}(\mathbb{C}) \cong \mathbb{C}^{\times}$ is a representation of $G$. Then $\rho(g) = \omega$ is a unit in $\mathbb{C}$ such that $\omega^{n} = 1$. Conversely, each $n$th root of unity $\omega$ determines a representation $\rho : G \to \mathbb{C}^{\times}$ by $g \mapsto \omega$.

The latter follows by the universal property of $G \cong \mathbb{Z} / n$ --- the quotient of the free group with one generator $g$ by the subgroup generated by $g^{n}$. Just as homomorphisms $\mathbb{Z} \to H$ (for any group $H$) are in bijection with elements of $H$, homomorphisms $\mathbb{Z} / n \to H$ are in bijection with elements $h \in H$ satisfying $h^{n} = 1$.

\newpage

%==========
\begin{tcolorbox}[title=Exercise 3]
Consider the usual action of $D_8$ on the four vertices of a square. Let $M$ be the associated $\mathbb{C}D_8$-permutation module. Find at least four distinct submodules of $M$. 
\end{tcolorbox}

\bigskip

For clarity, we write out the $\mathbb{C}[D_{8}]$-module structure on $M$ explicitly.

Let $e_{1}, \dots, e_{4}$ be a basis of $M \cong_{\mathsf{Mod}_{\mathbb{C}}} \mathbb{C}^{\oplus 4}$. Let $D_{8}$ act on it as follows
 \[
	\begin{tikzcd}
		e_{1} \ar[rr, "r"] \ar[rr, "s"', bend right] && e_{2} \ar[dd, "r"] \ar[ll, "s"', bend right] \\ \\
		e_{4} \ar[uu, "r"] \ar[rr, "s"', bend right] && e_{3} \ar[ll, "r"] \ar[ll, "s"', bend right]
	\end{tikzcd}
\]

Note, by Maschke's theorem it suffices to find two distinct submodules that are not complementary (our proof doesn't just find a complementary submodule, it finds a unique complementary submodule for each submodule $N \subseteq M$). The other two are their complements.

Consider $N_{1} \subseteq M$ spanned by $v = e_{1} + \dots + e_{4}$. Note $r \cdot v = v$ and $s \cdot v = v$ so this really is a $\mathbb{C}[D_{8}]$-submodule.

There is also $N_{2} \subseteq M$ spanned by $v_{1} = e_{1} + e_{3}$ and $v_{2} = e_{2} + e_{4}$. Note, $r \cdot v_{1} = s \cdot v_{1} = v_{2}$ and $r \cdot v_{2} = s \cdot v_{2} = v_{1}$. Thus, $N_{2}$ is a $\mathbb{C}[D_{8}]$-submodule.

Finally, since $N_{1}$ and $N_{2}$ have different $\mathbb{C}$-dimensions, they are distinct.


\newpage

%==========

\end{document}  
